\section{Proposed explainable algorithm}\label{sec:explainablealgorithm}
% Boran Tae 
        As discussed in the previous section, recent research has increasingly focused on explaining credit risk predictions. Techniques such as SHAP and LIME, commonly used for explaining credit risk predictions, assign importance scores to individual features; however, they do not account for combinations of feature values when explaining a prediction. This limitation highlights the need for approaches that not only assess the contribution of individual features when explaining predictions, but also combinations of them.\\

    {\color{blue}
		    	The idea of our algorithm is to integrate tree-based predictors with contrast pattern mining algorithms. Unlike traditional explainers for credit risk prediction, which operate independently of the predictor and rely on different techniques to generate explanations (as in techniques such as SHAP or LIME), in our algorithm the prediction and explanation processes are separated, but both are grounded in tree-based structures. In this way, although the explanation is not directly derived from the predictor, it preserves structural coherence with it, enabling the generation of explainations the structure of the data.\newline
		    	
		    	As mentioned previously, credit risk prediction depends on combinations of feature values rather than individual features considered in individually. For instance, a 'high debt-to-income ratio' may only indicate default when combined with 'low job stability'. While existing explanation techniques for credit risk prediction typically evaluate features individually, our algorithm represents combinations of feature values through contrast patterns. These patterns provide explanations based on conjunctions of feature conditions, similar to those used by tree-based predictors.\newline\newline
		    	
		    	Our choice of tree-based predictors due not only to their success in state-of-the-art credit risk prediction \cite{aruleba2025improved}, but also to their alignment with contrast patterns. This compatibility enables a transition between prediction and explanation. In particular, a decision tree can be interpreted as a hierarchy of decision rules, where each path from the root to a leaf node corresponds to a conjunction of conditions of the form $feature\#value$. Similarly, contrast patterns are expressed as conjunctions of conditions $feature\#value$.\newline
		    	
		    	By employing a contrast pattern mining algorithm compatible with this representation, a direct and consistent correspondence is maintained between the predictor and the generated explanations. %This structural alignment mitigates the limitations of approaches such as LIME or SHAP, which can produce explanations that do not fully capture the prediction.

		\subsection{First stage}
       	Because the datasets for credit risk prediction are imbalanced, and following the results reported in \cite{melchor2025oversampling}, we propose to apply SMOTENC \cite{chawla2002smote, mukherjee2021smote} to the training set to address class imbalance before training a predictor or extracting contrast patterns.\newline
       
       	For prediction, we propose using the FT4cip predictor \cite{canete2022ft4cip}, a tree-based predictor capable of handling mixed and imbalanced datasets, which has reported good results with this type of  dataset making it suitable for credit risk prediction. To construct the explanation, we propose using contrast patterns extracted from similar instances within each class. However, instances within the same class may exhibit different combinations of feature values, forming distinct subgroups within the class. Therefore, we propose extracting contrast patterns from clusters generated within each class of the dataset, allowing the identification of more specific and representative patterns for each subgroup.\newline
       	
       	To generate clusters, we first split the dataset into classes (defaulter and non-defaulter). Within each class, we propose using the Kmodes algorithm \cite{huang1997clustering, huang1998extensions}. This choice is motivated by the need to group instances with similar feature-value combinations, particularly in mixed-type data. To determine the number of clusters in each class, we use the Silhouette coefficient \cite{dinh2019estimating, kuswardana2025comparison}, ensuring that the resulting clusters capture well-defined groupings within each class. This clustering step enables the extraction of contrast patterns that reflect localized structures within each class, improving the quality and interpretability of the explanations. \newline
       
       	For mining contrast patterns, we use CPMining, introduced in the PBC4cip algorithm \cite{loyola2017pbc4cip}. This contrast pattern miner is supervised, meaning that it requires labeled data and explicitly exploits differences between classes to identify discriminative patterns. Thus, we propose applying it on each cluster of a given class against the other class, rather than using the entire class as a whole. This design choice is motivated by two main factors. First, by comparing each cluster against the opposite class, the extracted patterns remain class-discriminative, which is consistent with the supervised nature of the mining algorithm. Second, operating at the cluster level allows capturing more specific combinations of feature values that characterize subgroups within the class, instead of producing overly general patterns. As a result, the generated contrast patterns are both cluster-specific and class-discriminative, improving their representativeness and explainability for explanation purposes.\newline
       
       	After mining, we propose a filtering process to remove redundant contrast patterns. Following the filtering strategy used in CPMining, redundancy is identified by comparing pairs of patterns based on their conditions. Given two patterns $P_{1}$ and $P_{2}$, $P_{1}$ is considered more general than $P_{2}$ if its set of conditions is a subset of those defining $P_{2}$, in which case $P_{2}$ is more specific. Applied at the cluster level, this process effectively removes redundant patterns while preserving the most representative and informative ones, thereby enhancing the explainability of the extracted patterns and reducing their overall number.\newline
       
       	For example, consider the next contrast patterns extracted from the German Credit Data dataset \cite{statlog_german_credit_data_144}:\newline
       
       	$P_{1}:$ Credit Duration $\leq$ 15 months \textbf{and} Property$=$real estate \textbf{and} Foreign Worker$=$yes\newline
       
       	and the more specific contrast pattern\newline
       
       	$P_{2}:$  Credit Duration $\leq$ 15 months \textbf{and} Property$=$real estate \textbf{and} Foreign Worker$=$yes \textbf{and} Checking Account $\neq$ low\newline
       
       	In this case $P_{1}$ is more general than $P_{2}$. If both contrast patterns cover the same instances within the cluster, the second contrast pattern does not provide additional information, since it introduces an extra condition while describing the same subset of instances. Therefore, only the more general contrast pattern is retained. \newline
       	
       	Additionally, contrast patterns whose support in the opposite class is greater than zero are removed, ensuring that the final set contains only contrast patterns that characterize the cluster while distinguishing it from the other class.
        
       %\subsection{Second stage}
       %In the proposed algorithm, the prediction and explanation of a new instance are performed sequentially. Given a new instance $x$, the prediction stage predicts the credit risk class of $x$ (defaulter or non-defaulter). Once the class $c$ has been determined, only the clusters associated with that class are considered for explanation.\newline
       	
       %To identify the cluster most similar to a new instance $x$, we define a contrast pattern-based similarity score between $x$ and each cluster $C_{c,j}$ associated with the predicted class $c$. The similarity score is defined in Equation~\ref{eq:cluster_similarity}.
       	
       %\begin{equation}
       %		Sim(x, C_{c,j})
       %		=
       %		\frac{
       %			\sum\limits_{p \in \mathcal{P}{c,j}(x)} supp(p)
       %		}{
       %			\sum\limits{p \in \mathcal{P}_{c,j}} supp(p)
       %		},
       %		\label{eq:cluster_similarity}
       %	\end{equation}
       	
       %Let $\mathcal{P}{c,j}$ denote the set of contrast patterns extracted from cluster $C{c,j}$, and let $\mathcal{P}{c,j}(x) \subseteq \mathcal{P}{c,j}$ be the subset of patterns that cover the instance $x$. The term $supp(p)$ represents the support of pattern $p$ within cluster $C_{c,j}$.\newline
       	
       %The similarity score is calculated as the sum of the supports of all contrast patterns in cluster $C_{c,j}$ that cover the instance $x$, divided by the total support of all contrast patterns extracted from that cluster. The cluster with the highest similarity score is selected as the most similar to the instance. Finally, an explanation is generated by selecting a representative contrast pattern from the pattern set of the selected cluster.\newline
       	
       %Once the most similar cluster $C_{c,j}$ has been identified, the explanation for the prediction is derived exclusively from this cluster. Let $\mathcal{P}_{c,j}(x)$ denote the set of all contrast patterns extracted from $C_{c,j}$ that cover the instance $x$. From this subset, the algorithm selects a pattern $p$ with the highest support, as defined in Equation \ref{eq:pattern_selection}.
       	
       %\begin{equation}
       %		p^* = \arg\max_{p \in \mathcal{P}_{c,j^*}(x)} supp(p),
       %		\label{eq:pattern_selection}
       %	\end{equation}
       	
       %The selected pattern $p^*$ constitutes the final explanation of the prediction. This contrast pattern serves as the most representative and discriminative description of the instance within the selected cluster, as it maximizes support among all patterns covering the instance.
    }